\documentclass[11pt]{article}
\usepackage[margin=1in]{geometry}
\usepackage{amsmath,amssymb,amsthm,booktabs}
\usepackage[hidelinks]{hyperref}
\setlength{\emergencystretch}{2em}
\newtheorem{theorem}{Theorem}
\newtheorem{proposition}[theorem]{Proposition}
\newtheorem{corollary}[theorem]{Corollary}
\newcommand{\mr}{\operatorname{mr}}
\newcommand{\ar}{\operatorname{ar}}
\newcommand{\Wh}{\operatorname{Wh}}
\newcommand{\C}{\mathbb C}
\title{Actual peripheral curves and iterated Whitehead doubles of torus knots}
\author{Internal research working proof and source-derived corollary}
\date{16 November 2026}
\begin{document}
\maketitle

For coprime $2\le p<q$, the positive torus knot $T(p,q)$ has an
explicit curve of irreducible rank-$p$ KCH representations with
\[
 \lambda=(-1)^{p+1}\mu^{-(p-1)q}.
\]
It supplies actual augmentation points satisfying the iterable-branch
condition of Gao--Gu--Li--Wang \cite{gglw}. Consequently their
Corollary~4.5 gives
\begin{equation}\label{eq:sharp}
 \ar(\Wh^N(T(p,q)))=\mr(\Wh^N(T(p,q)))
 =b(\Wh^N(T(p,q)))=p\,2^N \qquad(N\ge0).
\end{equation}
Here $\Wh$ is the fixed untwisted Whitehead pattern, zero framing and
peripheral convention of that source. The equality includes
augmentation rank over $\C$ at $U=1$.

This is a corollary of existing results, including a November 2026
preprint, with an explicit companion-matrix construction supplied
below. It is not a priority claim, an independent certification of
that entire preprint, a solution of full MRC, or a closure of any of
the retained sixteen table targets. Cornwell already constructs the
torus peripheral branch \cite{cor14}; his later theorem yields a
broader sufficient criterion \cite{cor17}. The finite controls check
specific algebraic identities, rather than establish the uniform
conclusions by sampling.

\section{Imported results and a broader sufficient criterion}

A KCH representation here is an irreducible complex representation
for which a meridian has eigenvalues $(\mu,1,\ldots,1)$, with
$\mu\ne1$. Its dimension bounds meridional rank from below
\cite[Theorem~1.1]{cor14}. Augmentations with $\mu\ne1$ correspond to
these representations, with the same rank \cite[Theorem~1.2]{cor17}.

We use the following precise input from \cite[Definition~4.1,
Corollary~4.5]{gglw}: an actual rank-$d$ augmentation peripheral curve
in coordinates $(\mu,\eta)=(\mu,\lambda^2)$, having an ideal point
with
\[
 \nu(\mu)=-u<0,\qquad \nu(\eta)=v,
 \qquad v>u\ \hbox{or}\ v<0,
\]
gives augmentation rank at least $2^Nd$ for every finite $N$.
If $d=b(K)$, all three ranks of $\Wh^N(K)$ equal $2^Nd$.
Actual means that a dense subset consists of peripheral data of
rank-$d$ augmentations; a factor in a polynomial alone is insufficient.
The initial point may depend on $N$.

\begin{proposition}[A direct Cornwell--GGLW consequence]\label{prop:braid}
Let a nontrivial knot $K$ be the closure of an $n$-strand braid $B$.
Suppose $K$ has a rank-$n$ complex augmentation with $\mu\ne1$ and
the braid writhe $w(B)$ is nonzero. Then
\[
 \ar(\Wh^N(K))=\mr(\Wh^N(K))=b(\Wh^N(K))=n\,2^N
 \quad\text{for every }N\ge0,
\]
\end{proposition}
\begin{proof}
Cornwell's Theorem~1.10 \cite{cor17} gives the curve
$\lambda=(-\mu)^{-w(B)}$. Importantly, the last paragraph of its
proof constructs actual rank-$n$ augmentations for all but finitely
many $\mu\in\C^*$: the determinant controlling rank is a polynomial,
nonzero at the initial augmentation. Thus $\eta=\mu^{-2w(B)}$ is an
actual peripheral curve. At $\mu=\infty$ its orders are
$\nu(\mu)=-1$ and $\nu(\eta)=2w(B)$. Nonzero integral writhe gives
either $2w(B)>1$ or $2w(B)<0$.
Finally $n\le\mr(K)\le b(K)\le n$, using the augmentation bound
and the closed-braid upper bound. The imported iterable-curve result
now applies with $d=n$.
\end{proof}

\section{A uniform companion-matrix torus family}

The construction below realizes the maximal torus branch directly.
It avoids relying on an unlabelled factor of an augmentation
polynomial for the actual-rank hypothesis.

Use the standard positive torus presentation and preferred peripheral
elements \cite[Section~5.1]{cor14}:
\[
 \pi_{T(p,q)}=\langle x,y\mid x^p=y^q\rangle,
 \qquad m=x^a y^b,\quad aq+bp=1,\quad \ell=x^pm^{-pq}.
\]
For a monic polynomial $h$ of degree $p$, let $C_h$ be its companion
matrix with subdiagonal ones and last column the negatives of its
coefficients in ascending order. Choose distinct $q$th roots of
unity $\delta_1,\ldots,\delta_p$ and put
\[
 F=C_{z^p-1},\qquad
 C(s)=C_{\prod_j(z-s\delta_j)},\qquad D(s)=s^{-1}C(s)
 \quad(s\in\C^*).
\]
The matrices have simple spectra, so $F^p=I$ and $D(s)^q=I$.
Choose $r,t$ with $-ar\equiv1\pmod p$ and $bt\equiv1\pmod q$.
These choices exist because $a$ is a unit modulo $p$ and $b$ is a
unit modulo $q$. Define
\[
 X(s)=s^qF^r,\qquad Y(s)=s^pD(s)^t.
\]
Then $X^p=Y^q=s^{pq}I$, giving a representation, and
\begin{equation}\label{eq:meridian}
 M(s)=X^aY^b=sF^{-1}D(s)=F^{-1}C(s),\qquad
 M(s)-I=F^{-1}(C(s)-F).
\end{equation}
Since the first $p-1$ companion columns agree, the last difference
has rank at most one. Let
\[
 c=\frac{\det D(s)}{\det F}
   =(-1)^{p+1}\prod_j\delta_j.
\]
Away from $cs^p=1$, $M-I$ has rank exactly one and $M$ is
diagonalizable with nonidentity eigenvalue $\mu=cs^p$.

\begin{proposition}\label{prop:irred}
If the spectra of $F$ and $C(s)$ are disjoint and $cs^p\ne1$,
the representation $(X(s),Y(s))$ is an irreducible KCH representation
of dimension $p$.
\end{proposition}
\begin{proof}
Both conditions exclude only finitely many $s$. Invariance under
$X,Y$ is equivalent to invariance under $F,C(s)$: $r$ and $t$ are
units modulo their respective orders, and multiplication by nonzero
scalars does not change invariant subspaces.

Write $C(s)-F=u e_p^T$. A common invariant subspace contained in
$\ker e_p^T$ is zero, because the rows $e_p^TF^j$ span the dual
space. Any other common invariant subspace contains $u$. For each
$p$th root $\alpha$ of unity, the row
$v_\alpha^T=(1,\alpha,\ldots,\alpha^{p-1})$ is a left
$\alpha$-eigenvector of $F$. Direct companion multiplication gives
\[
 v_\alpha^Tu=(\alpha^p-1)-\prod_j(\alpha-s\delta_j)
 =-\prod_j(\alpha-s\delta_j)\ne0.
\]
As $F$ has distinct eigenvalues, this says $u$ has nonzero projection
to every eigenline of $F$. Hence $u$ is cyclic for $F$, and its
invariant span is the entire space. This proves irreducibility;
equation~\eqref{eq:meridian} proves the KCH condition.
\end{proof}

On the distinguished meridian line, the longitude eigenvalue is
\[
 \lambda=s^{pq}\mu^{-pq}.
\]
Consequently
\[
 \lambda\mu^{(p-1)q}=c^{-q}=(-1)^{p+1}.
\]
For the last sign, $c^q=(-1)^{(p+1)q}$; this equals
$(-1)^{p+1}$ because $p,q$ are coprime. As $s$ varies outside a
finite set, its image covers all but finitely many points of this
monomial peripheral curve. The rank is $p$ throughout that open set.

\begin{corollary}
Equation~\eqref{eq:sharp} holds, using the imported GGLW
iterable-curve theorem with its stated preprint status.
\end{corollary}
\begin{proof}
Here $\eta=\mu^{-2(p-1)q}$. The smooth projective model of this
curve has the point $\mu=\infty$, with orders
$\nu(\mu)=-1$ and $\nu(\eta)=2(p-1)q>1$. Proposition~\ref{prop:irred}
supplies the actual rank-$p$ points. The positive $p$-strand torus
braid gives $b\le p$, while this representation gives $\mr\ge p$;
hence $b=\mr=p$. The imported result applies. Alternatively,
Proposition~\ref{prop:braid} applies with $n=p$ and $w=(p-1)q$.
\end{proof}

\section{An actual varying-eigenvalue curve through an icosahedral point}

For $T(3,5)$, take $\tau^2=\tau+1$ and
\[
 F=\begin{pmatrix}0&0&1\\1&0&0\\0&1&0\end{pmatrix},\qquad
 C(s)=\begin{pmatrix}
 0&0&s^3\\1&0&(\tau-1)s^2\\0&1&(1-\tau)s
 \end{pmatrix},\qquad D=C(s)/s.
\]
The eigenvalues of $D$ are $1,\zeta_5^2,\zeta_5^3$ for the
corresponding choice of primitive root and embedding of $\tau$.
Set $X=s^5F$ and $Y=s^3D^3$. Since $m=x^{-1}y^2$,
\[
 M=F^2C(s)=\begin{pmatrix}
 1&0&(\tau-1)s^2\\0&1&(1-\tau)s\\0&0&s^3
 \end{pmatrix},\qquad
 \mu=s^3,\quad\lambda=s^{-30},\quad\eta=s^{-60}.
\]
At $s=-1$, $X^3=Y^5=-I$ and $M^2=I$. Exact enumeration over
$\mathbb Z[\tau]$ gives image order $120$. Modulo its central
$\{\pm I\}$, it is the $(3,5,2)$ icosahedral group, so the image
is the three-dimensional $H_3$ reflection image.
Both peripheral functions vary:
\[
 \left.\frac{d\mu}{ds}\right|_{s=-1}=3,
 \qquad \left.\frac{d\lambda}{ds}\right|_{s=-1}=30.
\]
This is an actual algebraic curve, not just an infinitesimal
direction. For generic $s$, $\mu^2\ne1$, and a rank-one meridian
with such an eigenvalue cannot preserve a nondegenerate symmetric
form. Thus this curve escapes the fixed-involution orthogonality
boundary of the companion note. It is a torus control, not a curve
for any member of the retained sixteen-knot panel.

\section{Checks, source assurance and remaining scope}

The frozen plan \path{torus_whitehead_plan.json} specifies thirteen
coprime torus pairs, with three finite-field samples each. Exact
matrix arithmetic checks the torus relation, meridian word,
rank-one condition, preferred longitude eigenvalue and full matrix
algebra dimension $p^2$. Symbolic identities over
$\mathbb Q(\tau)(s)$ verify the displayed $T(3,5)$ family, including
its $120$-element specialization.

For the source doubling block, separate symbolic polynomial-remainder
checks verify the boundary commutator and the longitude word on its
distinguished eigenvector. The frozen $2,976$ finite-field triples
over primes $5,7,11,13$ yield respectively $0,4,64,88$ admissible
blocks; all satisfy the checked identities. No field or target was
replaced after a miss. These are algebraic controls, with shared
symbolic-library infrastructure; they do not independently verify
the topology of the Whitehead peripheral identification.

The source audit read the complete one-step linear-algebra proof,
the peripheral calculation and the iterable-branch argument in
\cite[Sections~2.3--2.4 and 4.1]{gglw}, checking the latter's order
calculations and nonvanishing conditions. Its pattern presentation,
zero-framed longitude and geometric bridge construction remain
explicit imported topological inputs. The full microlocal theory,
two-bridge existence argument and all other parts of the 44-page
preprint are not independently certified. Both immutable HTML and
PDF snapshots are retained with hashes; their author-date lines
differ (11 and 13 November), without an inferred mathematical repair.

The gap in Dutra's 2022 Lemma~3.1 remains separately checked. This
corollary has no dependency on that lemma or the broader once-double
claim affected by it. No full MRC theorem or previously unresolved
table-knot closure is asserted. Priority beyond the cited direct
dependencies, formal verification and specialist review remain open.
The next existence gate is an actual sharp-rank peripheral curve
for a family outside the imported criteria, or for one of the
retained sixteen targets; the unexecuted tangent proposal supplies
no evidence of such a curve.

\begin{thebibliography}{9}
\bibitem{cor14} C. Cornwell, \emph{Knot contact homology and
representations of knot groups} (2014), Theorems~1.1, 1.3 and
Section~5.1. \url{https://arxiv.org/abs/1303.4943v5}.
\bibitem{cor17} C. Cornwell, \emph{KCH representations,
augmentations, and A-polynomials}, J. Symplectic Geometry 15
(2017), 983--1017, Theorems~1.2 and 1.10; arXiv proof p.~18.
\url{https://arxiv.org/abs/1310.7526v3}.
\bibitem{gglw} H. Gao, S. Gu, Z. Li and J. Wang,
\emph{Meridional ranks of iterated Whitehead doubles}, November 2026
preprint, Sections~2.3--2.4 and 4.1, especially Corollary~4.5.
\url{https://arxiv.org/abs/2611.15673v1}.
\end{thebibliography}
\end{document}
